%% HEADER
%%%%%%%%%%%%%%%%%%%%%%%%%%%%%%%%%%%%%%%%%%%%%%%%%%%%%%%%%%%%%
%%
%%
\newcommand{\hyperrefpdfauthor}{}
\newcommand{\hyperrefpdftitle}{}
\newcommand{\hyperrefpdfsubject}{}
\newcommand{\hyperrefpdfkeywords}{}
\newcommand{\hyperrefpdfborder}{0}
\newcommand{\subfigureautorefname}{\figureautorefname}
\documentclass{styles/wissdoc-kw-eng}

\usepackage{pdfpages}
\usepackage[printonlyused]{acronym}

% % newer acronym packaged dont use bflabel anymore,
\providecommand\aclabelfont{} % the new command ist aclabelfont however, we also need to be backwards compatible so... we do it like this
\renewcommand{\aclabelfont}[1]{\normalfont{\normalsize{#1}}\hfill} % keine serifenlose schrift für acronym

% newer acronym packaged dont use bflabel anymore, to not fail the statement below provide the command if not existing
\providecommand\bflabel{} 
\renewcommand{\bflabel}[1]{\aclabelfont{#1}} 

%\usepackage{acronym}
\usepackage{subfigure}

\usepackage{float}
\floatstyle{ruled}
\newfloat{listing}{htbp}{lop}[chapter]
\floatname{listing}{Listing}

\usepackage[hang,center,nooneline]{caption}
\captionsetup[figure]{font={small,sf}}
\captionsetup[table]{font={small,sf}}
\captionsetup[listing]{font={small,sf}}
\usepackage[absolute]{textpos}
\usepackage{styles/etoolbox}


%% Normales LaTeX oder pdfLaTeX? %%%%%%%%%%%%%%%%%%%%%%%%%%%%
%% ==> Das neue if-Kommando "\ifpdf" wird an einigen wenigen
%% ==> Stellen benötigt, um die Kompatibilität zwischen
%% ==> LaTeX und pdfLaTeX herzustellen.
%\newif\ifpdf
%\ifx\pdfoutput\undefined
%    \pdffalse              %%normales LaTeX wird ausgeführt
%\else
%    \pdfoutput=1
%    \pdftrue               %%pdfLaTeX wird ausgeführt
%\fi


%% Fonts für pdfLaTeX %%%%%%%%%%%%%%%%%%%%%%%%%%%%%%%%%%%%%%%
%% ==> Nur notwendig, falls keine cm-super-Fonts installiert
\ifpdf
	\usepackage{ae}       %%Benutzen Sie nur eines dieser Pakete:
	%\usepackage{zefonts}  %%je nachdem, welches Sie besitzen.
\else
	%%Normales LaTeX - keine speziellen Fontpackages notwendig
\fi


%% Deutsche Anpassungen %%%%%%%%%%%%%%%%%%%%%%%%%%%%%%%%%%%%%
%\usepackage[T1]{fontenc}
%\usepackage[utf8]{inputenc}

%% zur Zitaten des Quelltextes%%%%%%%%%%%%%%%%%%%%%%%%%%%%%%%
% "final" forces printing of all listings, even if the global "draft" is set
\usepackage[final]{listings}
\lstset{
    basicstyle=\footnotesize\ttfamily,
    tabsize=4,
    numberstyle=\tiny\color{gray},
    numbersep=5pt,
    numbers=left,
    captionpos=b,
    abovecaptionskip=0pt,
    belowcaptionskip=0pt,
    aboveskip=10pt,
    belowskip=0pt,
    floatplacement=tbp,
    frame=topline,
    framerule=.1pt,
    framesep = 3pt,
    }
\renewcommand\lstlistingname{\textbf{Listing}}
% This is only kept for backwards compatibility. You should never have to use it. Use the listing-environment instead.
%\DeclareCaptionFormat*{lstruled}{{\bfseries#1\small\space\normalfont#3\hrule height.1pt depth0pt}\par}
%\captionsetup[lstlisting]{format=lstruled,singlelinecheck=false}

%% mehrere Abbildungen in eine %%%%%%%%%%%%%%%%%%%%%%%%%%%%%%
\usepackage{subfigure}

%% Packages für Formeln %%%%%%%%%%%%%%%%%%%%%%%%%%%%%%%%%%%%%
\usepackage{amsmath}
\usepackage{amsthm}
\usepackage{amsfonts}

%% Zeilenabstand %%%%%%%%%%%%%%%%%%%%%%%%%%%%%%%%%%%%%%%%%%%%
\usepackage{setspace}
%\singlespacing        %% 1-zeilig (Standard)
%\onehalfspacing       %% 1,5-zeilig
%\doublespacing        %% 2-zeilig


%% Andere Packages %%%%%%%%%%%%%%%%%%%%%%%%%%%%%%%%%%%%%%%%%%
%\usepackage{a4wide} %%Kleinere Seitenränder = mehr Text pro Zeile.
\usepackage{fancyhdr} %%Fancy Kopf- und Fußzeilen
%\usepackage{longtable} %%Für Tabellen, die eine Seite überschreiten
\usepackage{lscape}
\usepackage{rotating} 
%\usepackage[htt]{hyphenat} %Trennung von Typewriter-Schriften
%\usepackage{listings}
%\usepackage{pstricks-add}
\usepackage{tikz}


% Tabellen mit Center und left
\usepackage{tabularx,colortbl} % colored table background
\newcolumntype{C}[1]{>{\centering\arraybackslash}p{#1}}
\newcolumntype{R}[1]{>{\raggedleft\arraybackslash}p{#1}}
% Table spacings
\newcommand\T{\rule{0pt}{2.5ex}\rule[-1.0ex]{0pt}{0pt}}
\newcommand\B{\rule[-1.0ex]{0pt}{0pt}}

\definecolor{slightgray}{gray}{.90} 


%% Definitionen %%%%%%%%%%%%%%%%%%%%%%%%%%%%
\newcommand{\todo}[1]{ \textcolor{red}{\textbf{TODO: #1}} }

%% zur Benutzung bei ergänzenden Daten%%%%%%%%%%%%%%%%%%%%%%%%
%\usepackage{endnotes}
%\renewcommand{\notesname}{Konfigurationsdaten der Messreihen}
%\renewcommand{\theendnote}{\Alph{endnote}}
%\renewcommand{\enotesize}{\normalsize}

%\hyphenation{Sensor-netz-werk
%}

%%%%%%%%%%%%%%%%%%%%%%%%%%%%%%%%%%%%%%%%%%%%%%%%%%%%%%%%%%%%%
%% DOKUMENT
%%%%%%%%%%%%%%%%%%%%%%%%%%%%%%%%%%%%%%%%%%%%%%%%%%%%%%%%%%%%%
\begin{document}

%% Dateiendungen für Grafiken %%%%%%%%%%%%%%%%%%%%%%%%%%%%%%%
%% ==> Sie können hiermit die Dateiendung einer Grafik weglassen.
%% ==> Aus "\includegraphics{titel.eps}" wird "\includegraphics{titel}".
%% ==> Wenn Sie nunmehr 2 inhaltsgleiche Grafiken "titel.eps" und
%% ==> "titel.pdf" erstellen, wird jeweils nur die Grafik eingebunden,
%% ==> die von ihrem Compiler verarbeitet werden kann.
%% ==> pdfLaTeX benutzt "titel.pdf". LaTeX benutzt "titel.eps".
%\ifpdf
%    \DeclareGraphicsExtensions{.pdf,.jpg,.png}
%\else
%    \DeclareGraphicsExtensions{.eps}
%\fi

\pagestyle{empty} %%Keine Kopf-/Fusszeilen auf den ersten Seiten.

\ifnotdraft{
%% Deckblatt %%%%%%%%%%%%%%%%%%%%%%%%%%%%%%%%%%%%%%%%%%%%%%%%
\frontmatter
% !TEX root =  ../thesis.tex

\titlehead{
%	\hfill
%	\includegraphics[width=10cm,keepaspectratio]{logos/rwth}
} % end titlehead


\begin{titlepage}

\begin{textblock*}{10cm} (12.7cm,2.7cm)
\includegraphics[width=10cm,keepaspectratio]{logos/rwth}
\end{textblock*}

\let\footnotesize\small \let\footnoterule\relax

\hbox{}
\vfill

\centering

\begin{doublespace} 
{ \huge\sffamily\textbf{Thesis title \\ \vspace{-0.7em}
goes here\\ \vspace{0.2em}
}}
\end{doublespace}
\vskip 2cm

{\large\sffamily

Bachelor Thesis\\[5pt]
\textbf{First and Last name}
\vskip 1cm

This work was submitted to the\\[5pt]
\textbf{Department of Medical Informatics\\[5pt]
        RWTH Aachen University, Germany}
\vskip 2cm

Adviser(s):
\vskip 2mm
Dr. Stephan Jonas\\
Dr. Ekaterina Kutafina
\vskip 5mm
Examiners:
\vskip 2mm
Univ.-Prof. Dr.med. Dr.rer.nat. Dipl.-Math. Klaus Kabino\\
Prof. Dr. Examiner Second
\vskip 1cm

\begin{tabular}{R{6cm}p{6cm}}
Registration date:  & 2016-11-18 \\
Submission date:    & 2017-03-09 \\
\end{tabular}

} %\large\sffamily

\vfill

\end{titlepage}

\cleardoublepage
% !TEX root =  ../thesis.tex
\includepdf[pages={1},offset=75 -75,pagecommand={
\begin{tikzpicture}[remember picture,overlay,shift={(current page.center)}]
\coordinate (begin-generic) at (5.3,11.3);
\coordinate (end-generic) at (6.4,11.3);
\coordinate (begin-bachelor) at (6.45,11.3);
\coordinate (end-bachelor) at (9,11.3);
\coordinate (begin-master) at (-5.4,10.8);
\coordinate (end-master) at    (-3.15,10.8);
\node at (-5.5,13.4)
       [text width=7cm,rounded corners,below right]
{
Last name, First Name
};
\node at (3.25,13.4)  [text width=5cm,rounded corners,below right]
{
322352 % This is optional
};
% ===== for type generic uncomment the next 2 lines =====
% \draw[line width=1mm] (begin-bachelor) -- (end-bachelor);
% \draw[line width=1mm] (begin-master) -- (end-master);
% ===== for type bachelor uncomment the next 2 lines =====
\draw[line width=1mm] (begin-generic) -- (end-generic);
\draw[line width=1mm] (begin-master) -- (end-master);
% ===== for type master uncomment the next 2 lines =====
% \draw[line width=1mm] (begin-generic) -- (end-generic);
% \draw[line width=1mm] (begin-bachelor) -- (end-bachelor);
%
\node at (-5.5,10.4)
       [text width=16.5cm,rounded corners,below right]
{
Thesis title goes here
};
\node at (-5.5,4.6)
       [text width=7cm,rounded corners,below right]
{
Aachen, 09.03.2017
};
\node at (-5.5,-8.45)
       [text width=7cm,rounded corners,below right]
{
Aachen, 09.03.2017
};
\end{tikzpicture}
}]{Formular_Eidesstattliche_Versicherung_neu.pdf}

\cleardoublepage
\cleardoublepage

\begin{center}
\paragraph{Kurzfassung}
\hrulefill
\end{center}
Elektroenzephalographie (EEG) ist eine Methode um verschiedene psychische Gesundheitsprobleme wie Demenz, Depression, Schizophrenie oder posttraumatische Belastungsstörung zu untersuchen....


\vspace {2cm}
\begin{center}
\paragraph{Abstract}
\hrulefill
\end{center}
Electroencephalography (EEG) is one method to investigate different mental health problems, like dementia, depression, schizophrenia or post-traumatic stress disorder...
\cleardoublepage
\cleardoublepage

\chapter*{Acknowledgments}

I want to thank Professor and Dr Bas, especially Bas. He encourages me all the time when the work turned stuck. He guided me through an wide topic to an clear solution. 

At last, I'm also thankful for the help from the colleges in PADS at RWTH Aachen. 
\cleardoublepage

% Titelseite hatte noch normale Tabellen. Von hier ab sollen alle
% Tabellen laut style-Vorgaben sans serif sein.
\AtBeginEnvironment{tabular}{\sffamily}
\AtBeginEnvironment{tabularx}{\sffamily}

%% Inhaltsverzeichnis %%%%%%%%%%%%%%%%%%%%%%%%%%%%%%%%%%%%%%%
\tableofcontents %Inhaltsverzeichnis
\cleardoublepage %Das erste Kapitel soll auf einer ungeraden Seite beginnen.
} % end ifnotdraft

\pagestyle{fancy} %%Ab hier die Kopf-/Fusszeilen: headings / fancy / ...

%%%%%%%%%%%%%%%%%%%%%%%%%%%%%%%%%%%%%%%%%%%%%%%%%%%%%%%%%%%%%
% einzelne Kapitel
%%%%%%%%%%%%%%%%%%%%%%%%%%%%%%%%%%%%%%%%%%%%%%%%%%%%%%%%%%%%%
%\input{commands}

\mainmatter
\chapter{Introduction}
Being the most complex organ in the body, the brain consists of roughly 10 billions of neurons\cite{mueller2000} ...
\\

\includegraphics{figures/rwth}
\chapter{Background}

\input{chapters/03_preliminary}




%%%%%%%%%%%%%%%%%%%%%%%%%%%%%%%%%%%%%%%%%%%%%%%%%%%%%%%%%%%%%
%% LITERATUR UND ANDERE VERZEICHNISSE
%%%%%%%%%%%%%%%%%%%%%%%%%%%%%%%%%%%%%%%%%%%%%%%%%%%%%%%%%%%%%
%% Ein kleiner Abstand zu den Kapiteln im Inhaltsverzeichnis (toc)
\ifnotdraft{
\addtocontents{toc}{\protect\vspace*{\baselineskip}}
\cleardoublepage
%% Literaturverzeichnis
\phantomsection % phantomsection wird benötigt, damit z.B. hyperref die richtige Seite verlinkt.
\addcontentsline{toc}{chapter}{Bibliography}
%\nocite{*} %Auch nicht-zitierte BibTeX-Einträge werden angezeigt.
\bibliography{literature/literature}%Eine Datei 'literatur.bib' wird hierfür benötigt.
\bibliographystyle{styles/vancouver}%Art der Ausgabe: plain / apalike / amsalpha / ...
}

%% Abbildungsverzeichnis
\clearpage
\addcontentsline{toc}{chapter}{List of Figures}
\listoffigures

%% Tabellenverzeichnis
\clearpage
\addcontentsline{toc}{chapter}{List of Tables}
\listoftables


%%%%%%%%%%%%%%%%%%%%%%%%%%%%%%%%%%%%%%%%%%%%%%%%%%%%%%%%%%%%%
%% ANHÄNGE
%%%%%%%%%%%%%%%%%%%%%%%%%%%%%%%%%%%%%%%%%%%%%%%%%%%%%%%%%%%%%
\appendix
\chapter{Appendix}

\section{List of Abbreviations}

\begin{acronym}[ABCDEFGH]
	\setlength{\itemsep}{-\parsep}
	
	\acro{EEG}{Electroencephalography}
	\acro{EDF}{European Data Format}

\end{acronym}

% \includepdf[pages=-,offset=-2.5cm -3cm ]{figures/appendix1.pdf}


\end{document}
